% ===================== 6. Discussion =====================
\section{Discussion}\label{sec:discussion}

\paragraph{Why is the trained gate inert here?}
The T2T gate of LLaDA2.2 was trained, via L-EBPO, in a regime where editing
happens \emph{inside} the current block while parallel denoising is still
resolving it~--- its rewrites correct early sampling noise within a block that is
still being written. Our schedule relocates those rewrites to a window behind a
one-token frontier, a distribution the gate was never trained on; the model's
confidences behind the frontier are apparently already well calibrated~--- the
gate fires (6.9 rewrites per task) but its rewrites are almost always
confidence-neutral polish, not error repair. The result echoes the diagnosis of
\citet{remeval}: the bottleneck is the dLLM's ability to \emph{detect} its own
mistakes, not the schedule through which revisions are applied. A gate trained
with a sliding window~--- or any trained mistake detector~--- is the natural next
experiment, and our kernel is ready to host it.

\paragraph{What the negative result buys.}
First, a clean baseline: any future claim that sliding-window revision helps a
LLaDA2-class model must now beat an honestly-paired null, and our segment analysis
(Section~\ref{sec:artifact}) gives reviewers a concrete failure mode~--- unequal
arm coverage~--- to check for. Second, a cost bound: the revision opportunity is
available at a predictable $\times 4.88$ NFE, and a $k$-token frontier would
divide that multiplier by $k$ without touching a single formula; if a future
trained gate does help, the engineering path to it is already cheap. Third, the
zero-invention methodology itself: every component of our decoder is verbatim
model code with runtime equivalence checks, which is what made the null result
interpretable rather than confounded by an accidental model change.

% ===================== 7. Limitations =====================
\section{Limitations}\label{sec:limitations}

Our conclusions are bounded by four constraints. (i) \emph{Sample size}: the main
arm has $n{=}45$ and the control $n{=}22$; Wilson intervals span $\pm 12$ points,
so we can exclude large effects but not effects of a few points; detecting a
$\pm 5$-point effect would require $n \geq 500$ per arm. (ii) \emph{One model,
one benchmark, one regime}: results are for LLaDA2.2-mini at 4-bit quantization
on GSM8K with greedy decoding and a 160-token cap that truncates most
generations; the HumanEval arm was skipped by the time budget. (iii)
\emph{Window sizes}: only $W{=}16$ (main) and $W{=}8$ (ablation, $n{=}10$) ran;
$W{=}32$ was not reached. (iv) \emph{Ablation identity}: the frontier-only
control differs from the wet-ink arm by the editing gate alone, but the wet-ink
family as a whole differs from stock decoding by the schedule; we disentangle
schedule cost (NFE) from editing effect (outcome flips), not the joint
interaction with sampling temperature, which we fixed at zero.

% ===================== 8. Conclusion =====================
\section{Conclusion}\label{sec:conclusion}

We implemented wet-ink decoding~--- a one-token left-to-right frontier with a
sliding window of revisable tokens and a hard dry edge~--- for a modern diffusion
language model, under a zero-invention policy in which every mechanism is the
base model's own trained code. Under a paired exact protocol the method preserves
the stock decoder's accuracy arm for arm (McNemar $p{=}1.0$ everywhere), its
apparent advantage over an internal control dissolves into a task-composition
artifact, and its cost is a predictable $\times 4.88$ NFE. Together with the four
earlier versions of this program, the evidence indicates that inference-time
token revision, by any rule so far published, does not improve a fixed diffusion
language model in the greedy, short-block regime; revision that helps will likely
have to be trained. Code, data, full telemetry, and the reproduction kernel are
public.\footnote{\url{https://github.com/Asilbekov/Explore} and
\url{https://huggingface.co/Freedomform/wet-ink-v5}}
